% !TeX program = xelatex
% Compile twice: xelatex triangular_commutator_Lp.tex
\documentclass[UTF8,11pt,a4paper,fontset=fandol]{ctexart}
\usepackage[margin=26mm]{geometry}
\usepackage{amsmath,amssymb,amsthm,mathtools}
\usepackage{xcolor}
\usepackage[colorlinks=true,linkcolor=blue!60!black,citecolor=blue!60!black,urlcolor=blue!60!black]{hyperref}
\usepackage{enumitem}
\setlist[enumerate]{itemsep=3pt,topsep=4pt}
\setlength{\emergencystretch}{2em}
\newtheorem{theorem}{定理}[section]
\newtheorem{lemma}[theorem]{引理}
\newtheorem{corollary}[theorem]{推论}
\theoremstyle{remark}
\newtheorem{remark}[theorem]{说明}
\newtheorem{example}[theorem]{例}
\renewcommand{\proofname}{证明}
\newcommand{\BH}{B(H)}
\newcommand{\Sp}{S_p(H)}
\newcommand{\Tr}{\operatorname{Tr}}
\newcommand{\dist}{\operatorname{dist}}
\newcommand{\PP}{\mathcal P}
\newcommand{\QQ}{\mathcal Q}
\newcommand{\DD}{\mathcal D}
\newcommand{\HH}{\mathcal H}
\newcommand{\Alg}{\mathcal A}
\newcommand{\CC}{\mathcal C}
\newcommand{\norm}[1]{\left\lVert #1\right\rVert}
\DeclareMathOperator{\rank}{rank}
\title{三角投影交换子在非交换 $L_p$ 上的有界性\\
\large 关于 Power (1985) 定理 2.1 的补充}
\author{}
\date{}

\begin{document}
\maketitle
\vspace{-2em}

\noindent\textbf{直接回答。}
在原文定理 2.1 的框架下，$X$ 已假定为 $H$ 上的有界算子。
若只要求交换子在 $1<p<\infty$ 的 $L_p$ 上有界，\textbf{不需要对 $X$ 增加任何条件}。
这里的 $L_p$ 是
\[
 L_p\bigl(B(H),\Tr\bigr)=S_p(H),
\]
即 Schatten 类，并非未指定测度空间的交换型函数空间 $L_p(\Omega)$。
下文先证明这一结论，再给出允许 $X$ 为形式块矩阵时的充要条件。

\section{原文的设置和直接结论}

设 $H$ 为可分复 Hilbert 空间，
\[
 0=P_0<P_1<P_2<\cdots,\qquad \rank P_n<\infty,
 \qquad P_n\longrightarrow I\quad\text{强算子收敛},
\]
并令 $E_n=P_n-P_{n-1}$。原文的巢代数为
\[
 \Alg=\{T\in\BH:(I-P_n)TP_n=0\ \text{对所有 }n\},
\]
即相对于 $H=\bigoplus_{n\geq1}E_nH$ 的块上三角代数。
在有限支撑块矩阵上定义
\[
 \PP(A)=\sum_{i\leq j}E_iAE_j,\qquad
 \QQ=I-\PP,\qquad
 \DD(A)=\sum_jE_jAE_j.
\]
当 $p=2$ 时，$\PP$ 就是原文投影到 $\Alg\cap S_2$ 的正交投影。
对每个 $1<p<\infty$，三角截断定理给出其唯一有界延拓，并且
\begin{equation}\label{eq:tp}
 t_p:=\norm{\PP}_{S_p\to S_p}<\infty,\qquad
 q_p:=\norm{\QQ}_{S_p\to S_p}<\infty.
\end{equation}
这些常数可以只依赖于 $p$，与块的个数和维数无关。
为使证明完整，附录给出这一事实的一个证明。

\medskip
原文 $X\PP-\PP X$ 中的 $X$ 表示\emph{左乘算子}
$L_X(A)=XA$。因此本文研究的是
\begin{equation}\label{eq:defC}
 \CC_X=[L_X,\PP],\qquad
 \CC_X(A)=X\PP(A)-\PP(XA),
\end{equation}
它与内部导子 $A\mapsto XA-AX$ 是不同的映射。

\begin{theorem}[保留原文的有界符号假设]\label{thm:bounded}
设 $X\in\BH$。则对每个 $1<p<\infty$，
$\CC_X$ 是 $S_p(H)$ 上的有界线性算子，且
\begin{equation}\label{eq:easybound}
 \norm{\CC_X}_{S_p\to S_p}\leq 2t_p\norm{X}_{\BH}.
\end{equation}
特别地，不要求 $X$ 紧、不要求 $X\in S_p$，也不要求
$X\in(\Alg+K(H))\cap(\Alg+K(H))^*$。
\end{theorem}

\begin{proof}
Schatten 理想不等式给出
\[
 \norm{XA}_p\leq\norm{X}_{\BH}\norm{A}_p,
 \qquad A\in S_p(H).
\]
结合 \eqref{eq:tp}，有
\begin{align*}
 \norm{X\PP(A)-\PP(XA)}_p
 &\leq\norm{X}_{\BH}\norm{\PP(A)}_p
       +t_p\norm{XA}_p\\
 &\leq 2t_p\norm{X}_{\BH}\norm{A}_p.
\end{align*}
这证明了结论。原文定理 2.1(ii) 中的代数条件刻画的是
$S_2$ 上的\emph{紧性}，而上述结论只涉及有界性。
\end{proof}

\section{若不预设 \texorpdfstring{$X$}{X} 有界：精确的符号条件}

本节的强化结论适用于原文的同一离散巢，包括有限维块的重数情形。
允许 $X=(X_{ij})_{i,j\geq1}$ 仅为形式块矩阵，
其中 $X_{ij}:E_jH\to E_iH$；此时不预设 $X$ 在 $H$ 上定义有界算子。
令 $X^{(N)}=(X_{ij})_{1\leq i,j\leq N}$，视为 $P_NH$ 上的算子，并定义
\begin{align}
 \beta_-(X)&=\sup_{N>n\geq1}
     \norm{(P_N-P_n)X^{(N)}P_n},\label{eq:beta-}\\
 \beta_+(X)&=\sup_{N>n\geq1}
     \norm{P_nX^{(N)}(P_N-P_n)},\label{eq:beta+}\\
 \beta(X)&=\max\{\beta_-(X),\beta_+(X)\}.\label{eq:beta}
\end{align}
这里所有范数均为 Hilbert 空间算子范数。
这些有限截断的定义避免了无界矩阵乘法的定义域问题。
若两个方向的角块已定义为有界算子，则等价地
\begin{equation}\label{eq:cut}
 \beta(X)=\sup_n\max\{
 \norm{(I-P_n)XP_n},\norm{P_nX(I-P_n)}\}
 =\sup_n\norm{[X,P_n]}.
\end{equation}
最后一个等号来自相对于 $P_nH\oplus(I-P_n)H$ 的表示
\[
 [X,P_n]=\begin{pmatrix}
 0&-P_nX(I-P_n)\\
 (I-P_n)XP_n&0
 \end{pmatrix}.
\]

对有限支撑 $A=(A_{kj})$，交换子的形式定义为
\begin{equation}\label{eq:entries}
 (\CC_X A)_{ij}
 =\sum_k X_{ik}A_{kj}
 \bigl(\mathbf1_{\{k\leq j\}}-\mathbf1_{\{i\leq j\}}\bigr).
\end{equation}
每个块的求和有限，但所得整个矩阵未必属于 $S_p$。
所谓其在 $S_p$ 上有界，是指存在 $S_p\to S_p$ 的有界映射，
在有限支撑矩阵上逐块满足 \eqref{eq:entries}。

\begin{theorem}[一致角块条件]\label{thm:sharp}
对上述形式块矩阵 $X$，下列条件等价：
\begin{enumerate}[label=\textup{(\roman*)}]
 \item 对某个 $1<p<\infty$，$\CC_X$ 在 $S_p(H)$ 上有界；
 \item $\beta(X)<\infty$；
 \item 对所有 $1<p<\infty$，$\CC_X$ 在 $S_p(H)$ 上有界。
\end{enumerate}
而且对每个 $1<p<\infty$，有
\begin{equation}\label{eq:sharpbounds}
 \boxed{\quad
 \beta(X)\leq\norm{\CC_X}_{S_p\to S_p}
 \leq t_pq_p\bigl(\beta_-(X)+\beta_+(X)\bigr)
 \leq2t_pq_p\beta(X).
 \quad}
\end{equation}
在 $p=2$ 时，还有精确等式
\begin{equation}\label{eq:exact2}
 \norm{\CC_X}_{S_2\to S_2}=\beta(X).
\end{equation}
\end{theorem}

先列出证明所需的有界完成引理。它使用原文在定理 2.1
之前已证明的距离公式
\begin{equation}\label{eq:nehari}
 \norm{\QQ L_Y\PP}_{S_2\to S_2}=\dist(Y,\Alg),
 \qquad Y\in\BH;
\end{equation}
见 \cite[p.~277]{Power}，其依据为原文定理 1.6 及其后的说明。
这里的距离取 $B(H)$ 算子范数。

\begin{lemma}[严格三角部分的有界完成]\label{lem:completion}
若 $\beta_-(X)<\infty$，则存在 $B\in\BH$ 满足
\[
 E_iBE_j=X_{ij}\quad(i>j),\qquad
 \norm B\leq\beta_-(X).
\]
若 $\beta_+(X)<\infty$，则存在 $C\in\BH$ 满足
\[
 E_iCE_j=X_{ij}\quad(i<j),\qquad
 \norm C\leq\beta_+(X).
\]
注意这里是为给定的严格三角部分\emph{补齐其他块}，
并未断言将其余块置零后所得的三角矩阵本身有界。
\end{lemma}

\begin{proof}
先说明对任意 $Y\in\BH$，
\begin{equation}\label{eq:Hcorner}
 \norm{\QQ L_Y\PP}_{S_2\to S_2}
 =\sup_j\norm{(I-P_j)YP_j}.
\end{equation}
事实上，逐块列计算得到
\[
 (\QQ L_Y\PP)(A)E_j=(I-P_j)YP_jAE_j.
\]
由 $\norm A_2^2=\sum_j\norm{AE_j}_2^2$ 得到 ``$\leq$''。
固定 $j$，取 $A=\xi\otimes\eta^*$，其中 $\xi\in P_jH$、
$\eta\in E_jH$ 均为单位向量，再对 $\xi$ 取上确界，得到反向不等式。
这里 $\xi\otimes\eta^*$ 表示 $h\mapsto\langle h,\eta\rangle\xi$。

将 $Y_N=X^{(N)}$ 在 $(I-P_N)H$ 上补零。
由 \eqref{eq:Hcorner}、\eqref{eq:nehari} 以及 \eqref{eq:beta-}，有
\[
 \dist(Y_N,\Alg)\leq\beta_-(X).
\]
故可取 $T_N\in\Alg$ 使
\[
 \norm{Y_N-T_N}\leq\beta_-(X)+N^{-1}.
\]
令 $B_N=P_N(Y_N-T_N)P_N$。则 $\norm{B_N}\leq\beta_-(X)+N^{-1}$，
且当 $N\geq i>j$ 时，$E_iB_NE_j=X_{ij}$。
由于 $H$ 可分，有界序列 $(B_N)$ 有弱算子收敛子列。
设其极限为 $B$，便有 $\norm B\leq\beta_-(X)$，
且每个严格下三角块均保持为 $X_{ij}$。

对形式伴随矩阵 $X^*=(X_{ji}^*)$ 应用刚证明的结论，
并使用 $\beta_-(X^*)=\beta_+(X)$，再取伴随，即得到 $C$。
\end{proof}

\begin{proof}[定理 \ref{thm:sharp} 的证明]
\textbf{必要性。}
设 $\CC_X$ 在某个 $S_p$ 上有界，记其范数为 $M$。
固定 $N>n$，取单位向量 $\eta\in E_nH$ 及任意 $\xi\in P_NH$，
令 $A=\xi\otimes\eta^*$。由 \eqref{eq:entries} 可得
\[
 P_N\CC_X(A)P_N
 =\bigl([X^{(N)},P_n]\xi\bigr)\otimes\eta^*.
\]
秩一算子的每个 Schatten 范数均等于两向量范数之积，
而压缩在 $S_p$ 上是压缩映射，故
\[
 \norm{[X^{(N)},P_n]\xi}
 \leq\norm{\CC_X(A)}_p\leq M\norm\xi.
\]
于是 $\norm{[X^{(N)},P_n]}\leq M$。
该范数正是两个相反角块范数的最大值。
对 $N,n$ 取上确界便得 $\beta(X)\leq M$。

\medskip
\textbf{充分性。}
设 $\beta(X)<\infty$。由引理 \ref{lem:completion} 取有界完成 $B,C$。
定义 $S_p$ 上的有界映射
\begin{equation}\label{eq:factorization}
 T=\QQ L_B\PP-\PP L_C\QQ.
\end{equation}
对有限支撑矩阵，直接计算得到
\begin{align*}
 (\QQ L_B\PP A)_{ij}
 &=\begin{cases}
     \displaystyle\sum_{k\leq j}X_{ik}A_{kj},&i>j,\\
     0,&i\leq j,
   \end{cases}\\
 (\PP L_C\QQ A)_{ij}
 &=\begin{cases}
     \displaystyle\sum_{k>j}X_{ik}A_{kj},&i\leq j,\\
     0,&i>j.
   \end{cases}
\end{align*}
其中第一式仅用到 $i>j\geq k$ 时 $B_{ik}=X_{ik}$，
第二式仅用到 $i\leq j<k$ 时 $C_{ik}=X_{ik}$。
因此 $T$ 逐块等于 \eqref{eq:entries}，从而就是所需交换子的延拓。
有限支撑矩阵在 $S_p$ 中稠密，故延拓唯一。
此外，
\[
 \norm T_{S_p\to S_p}
 \leq q_pt_p(\norm B+\norm C)
 \leq t_pq_p\bigl(\beta_-(X)+\beta_+(X)\bigr).
\]
这对每个 $1<p<\infty$ 均成立，证明了等价性和 \eqref{eq:sharpbounds}。

\medskip
\textbf{$S_2$ 上的等式。}
一致有界的角块定义了有界算子 $K_j=[X,P_j]$，且
$\norm{K_j}\leq\beta(X)$。
逐块列有
\[
 \CC_X(A)E_j=K_jAE_j.
\]
所以
\[
 \norm{\CC_X(A)}_2^2
 =\sum_j\norm{K_jAE_j}_2^2
 \leq\beta(X)^2\sum_j\norm{AE_j}_2^2
 =\beta(X)^2\norm A_2^2.
\]
结合已证的下界即得 \eqref{eq:exact2}。
\end{proof}

\begin{remark}[条件的含义及其与原文定理 2.3 的关系]
\leavevmode
\begin{enumerate}[label=\textup{(\arabic*)}]
 \item $\beta$ 是一个忽略块对角部分的半范数。
 若 $D=(D_{ij})$ 为形式块对角矩阵，即 $D_{ij}=0$ 对 $i\ne j$，则
 \[
  \CC_{X+D}=\CC_X,\qquad\beta(X+D)=\beta(X).
 \]
 因而允许形式符号时，不能把 ``$X\in B(H)$'' 当作必要条件。
 \item 在原文定理 2.3 所考虑的重数为一情形，将对角线归零以后，
 $\beta(X)<\infty$ 等价于
 \[
 X-\DD X\in B(H)+\widetilde{B(H)},
 \]
 其中 $\widetilde Y=-i\HH Y$ 是原文的三角共轭，
 $(\HH Y)_{ij}=\operatorname{sgn}(j-i)Y_{ij}$。
 这正是原文所讨论的 BMO 型符号条件；不能在未指定定义时
 将它直接替换为其他非交换 BMO 空间。
 \item 本文的充要条件针对原文的离散有限秩巢和左乘交换子。
 它没有把问题改成一般半有限 von Neumann 代数上的任意投影问题。
\end{enumerate}
\end{remark}

\section{两个区分有界性与紧性的例子}

\begin{example}[有界交换子可以不紧]
在 $H=\ell_2(\mathbb N)$ 的标准重数为一巢上，令
$Se_n=e_{n+1}$ 为单边移位。由于 $S\in B(H)$，定理 \ref{thm:bounded}
给出 $\CC_S$ 在所有 $1<p<\infty$ 的 $S_p$ 上有界。
但对矩阵单位 $e_{nn}$，有
\[
 \CC_S(e_{nn})=e_{n+1,n}.
\]
这些输入和输出的 $S_p$ 范数均为 $1$，且不同输出之间的距离为
$2^{1/p}$，故 $\CC_S$ 不紧。
\end{example}

\begin{example}[无界对角符号也可以产生有界交换子]
仍在标准巢上，取稠密定义的对角算子 $De_n=ne_n$。
其有限支撑矩阵交换子恒为零，故唯一延拓是 $S_p$ 上的零映射。
这说明精确的符号条件必须忽略对角部分。
此例按有限支撑定义交换子，不主张两个含无界乘子的项
分别对每个 $A\in S_p$ 都有定义。
\end{example}

\appendix
\section{三角截断在 \texorpdfstring{$S_p$}{Sp} 上有界的证明}

这里给出 \eqref{eq:tp} 的证明，所用基础工具只有 Schatten H\"older
不等式、对偶性和复插值。先在任意有限个有限维块上工作，
所有估计均与块的数目及维数无关。
设
\[
 (\HH A)_{ij}=\operatorname{sgn}(j-i)A_{ij},\qquad
 \operatorname{sgn}(0)=0.
\]
块对角映射 $\DD$ 在每个 $S_r$（$1\leq r\leq\infty$）上均为压缩映射：
它是块对角标量酉算子的共轭作用的平均。
此外，$\norm{\HH}_{S_2\to S_2}\leq1$。

\medskip
对任意块矩阵 $A,B$，逐项核对符号可得 Cotlar 型恒等式
\begin{equation}\label{eq:cotlar}
 (\HH A)(\HH B)
 =\HH\bigl((\HH A)B+A(\HH B)\bigr)-AB+(\DD A)(\DD B).
\end{equation}
其依据是对任意 $i,k,j$ 有
\begin{align*}
 &\operatorname{sgn}(k-i)\operatorname{sgn}(j-k)\\
 &\quad=\operatorname{sgn}(j-i)
   \bigl(\operatorname{sgn}(k-i)+\operatorname{sgn}(j-k)\bigr)
    -1+\mathbf1_{\{i=k=j\}}.
\end{align*}
由于 $\HH(A^*)=-(\HH A)^*$，在 \eqref{eq:cotlar} 中以 $A^*,A$
替代 $A,B$ 得到
\begin{equation}\label{eq:square}
 (\HH A)^*(\HH A)
 =A^*A-(\DD A)^*(\DD A)
   +\HH\bigl((\HH A)^*A-A^*\HH A\bigr).
\end{equation}
设已经知道 $\norm{\HH Z}_r\leq h_r\norm Z_r$，其中 $h_r$ 与维数无关。
对 \eqref{eq:square} 使用 $S_r$ 范数及 H\"older 不等式，可得
\begin{align*}
 \norm{\HH A}_{2r}^{,2}
 &\leq 2\norm A_{2r}^{,2}
       +2h_r\norm{\HH A}_{2r}\norm A_{2r}.
\end{align*}
解这一关于 $\norm{\HH A}_{2r}/\norm A_{2r}$ 的二次不等式，得到
\[
 \norm{\HH A}_{2r}\leq
 \bigl(h_r+\sqrt{h_r^2+2}\bigr)\norm A_{2r}.
\]
从 $h_2\leq1$ 出发归纳，便得到 $r=2,4,8,\ldots$ 时的统一有界性。
Schatten 类的复插值给出所有 $2\leq r<\infty$ 的有界性。
在 Hilbert--Schmidt 对偶配对下，$\HH$ 的伴随仍为 $\HH$，
因此对偶性又给出所有 $1<r<2$ 的有界性。

最后，
\[
 \PP=\tfrac12(I+\DD+\HH),\qquad
 \QQ=\tfrac12(I-\DD-\HH),
\]
故 $\PP,\QQ$ 在每个 $1<p<\infty$ 的 $S_p$ 上均一致有界。
利用有限支撑矩阵的稠密性即可延拓到整个 $S_p(H)$。
这就证明了 \eqref{eq:tp}，并完成定理 \ref{thm:bounded} 和
\ref{thm:sharp} 中所用的三角截断有界性论证。

\begin{thebibliography}{9}
\bibitem{Power}
S. Power,
\emph{Commutators with the triangular projection and Hankel forms on nest algebras},
J. London Math. Soc. (2) \textbf{32} (1985), 272--282.
\href{https://doi.org/10.1112/jlms/s2-32.2.272}{doi:10.1112/jlms/s2-32.2.272}.
本文使用其定理 2.1 之前的距离公式（第 277 页），并与定理 2.1、2.3 比较。
\end{thebibliography}

\end{document}
